\documentclass{article}
\usepackage[utf8]{inputenc}
\usepackage{amsmath}
\usepackage{geometry}
\geometry{margin=2cm}
\begin{document}

\section*{Análise da Questão 180 -- ENEM 2024 -- Matemática}

Um hospital possui 7 médicos cardiologistas e 6 médicos neurologistas. A direção quer formar uma equipe de 5 médicos para determinada atividade, sendo exigido que haja pelo menos 3 cardiologistas na equipe.

\subsection*{Solução Técnica}
Para atender à restrição de pelo menos 3 cardiologistas, consideramos as possíveis quantidades de cardiologistas que podem compor a equipe: 3, 4 ou 5.

\begin{itemize}
  \item Total de cardiologistas: 7
  \item Total de neurologistas: 6
  \item Tamanho da equipe: 5 médicos
\end{itemize}

Assim, o número total de equipes distintas é dado por:

\[
\binom{7}{3} \cdot \binom{6}{2} + \binom{7}{4} \cdot \binom{6}{1} + \binom{7}{5} \cdot \binom{6}{0}
\]

onde \(\binom{n}{k} = \frac{n!}{k!(n-k)!}\) é a combinação que calcula o número de formas de escolher \(k\) elementos entre \(n\).

\subsection*{Explicação Didática}
O enunciado exige que a equipe tenha pelo menos 3 cardiologistas, portanto precisamos somar todas as possibilidades que tenham de 3 até 5 cardiologistas.

\begin{enumerate}
  \item Quantidade total de cardiologistas e neurologistas no hospital.
  \item Dividir o problema considerando 3, 4 e 5 cardiologistas na equipe.
  \item Para cada caso, calcular a combinação de cardiologistas e de neurologistas necessárias para completar os 5 membros.
  \item Somar os resultados dos casos para obter todas as possibilidades.
\end{enumerate}

\subsection*{Erros Comuns}
\begin{itemize}
  \item Considerar só uma quantidade exata de cardiologistas, por exemplo apenas 3, ignorando o total de possibilidades.
  \item Confundir combinação com permutação, aplicando fórmulas erradas.
  \item Somar e multiplicar as combinações de forma incorreta ou errar nas quantidades escolhidas.
  \item Ignorar o termo "pelo menos", aplicando apenas uma condição fixa.
\end{itemize}

\subsection*{Dicas para Ensino}
Ensinar com exemplos pequenos e recursos visuais, além de reforçar a distinção entre combinação e permutação, ajuda a fixar o conceito e a interpretar corretamente "pelo menos" como soma de casos.

\subsection*{Análise dos Distratores}
As alternativas apresentadas são:

\begin{itemize}
  \item[A)] Produto de arranjos/fatoriais não correspondentes à soma dos casos.
  \item[B)] Multiplicação de combinações para a quantidade 3 cardiologistas e 2 neurologistas, sem somar outras possibilidades.
  \item[C)] Soma da fórmula das combinações para as três situações, porém com erro na estrutura ou nos símbolos.
\end{itemize}

A alternativa que melhor representa a solução correta, ou se aproxima dela, é a alternativa C.

\subsection*{Perfil da Questão}
Questão de nível médio voltada para combinatória aplicada, exige interpretação de condições e soma de casos para obter o número total de possibilidades.

\vspace{1em}

\textbf{Resposta correta (aproximada nas alternativas apresentadas):} C

\end{document}